\subsection{Typed terminal transfer via quotient-transition closure}

The terminal one-sort construction cannot be transferred merely by carrying
sort indices through the notation.  Cross-sort pair data can force a same-sort
quotient coordinate, and a typed word may return into a same-sort equivalence
without any individual cross-sort factor lying in that equivalence.  We now
build the missing finite transition closure explicitly.

For the hidden/HKN part of this subsection we harmlessly pass once more to the
finite quantifier-free definitional expansion naming every realised canonical
old tuple type on compatible real sorts up to the maximum relational arity.
This does not change embeddings, partial automorphisms, SAP, Ramsey arrows or
EPPA after reduct.  Its only purpose is that, in a completed visible
structure, the canonical old tuple-type relations partition every compatible
real tuple.  We suppress this definitional completion from the notation.

For current unary sorts \(a,b\), let \(\Rch^{ab}\) denote the finite family of
non-empty typed reachable profiles from \(a\) to \(b\).  Let
\(\Rel^{ab}\) denote the finite family of all non-empty \(0\)-definable
relations between these two transitive sorts, equivalently all non-empty
unions of realised canonical \(a\)-\(b\) pair types.  Every member of
\(\Rel^{ab}\), and in particular every reachable profile and every non-empty
exact domain, is total and surjective between its endpoint sorts.

\begin{lemma}[RAW saturation versus STRONG diversity equivalence factors]
\label{lem:typed-raw-strong-saturation}
Let $M$ be a homogeneous complete binary structure with individually
transitive named sorts, symmetric same-sort atomic relations and SAP.
No CCO or terminality hypothesis is needed for this lemma.
Let \(J,K\) be NONZERO actual
invariant same-sort equivalences on
named sorts \(a,b\), and put
\(J^\circ=J\setminus\Id_a\),
\(K^\circ=K\setminus\Id_b\).
For every actual old DIVERSITY
atomic relation \(t:a\to b\),
\[
       (J;\{t\};K)^\circ
        =J^\circ\starprod\{t\}
                          \starprod K^\circ.
\tag{TSC}\label{eq:typed-raw-strong-saturation}
\]
In addition, for two NONZERO actual
equivalences \(J,K\) on the SAME named
sort one has
\[
     (J;K)^\circ
           =J^\circ\starprod K^\circ
           =(J\vee K)^\circ.
\tag{TSC2}\label{eq:typed-equivalence-join-strong}
\]
This two-factor identity does NOT
introduce an artificial identity
as an intermediate STRONG atom.

With an identity equivalence at
either endpoint the corresponding
STRONG exterior factor is OMITTED.
Accordingly, in the following
transition-enriched typed subsection,
any nonzero equivalence \(J\) used
as a factor of \(\starprod\) means
its diversity palette \(J^\circ\);
in RAW inclusions and quotient
transport \(J\) continues to mean
the full relation including
logical identity. This is an
explicit notational coercion,
NOT an extra identity atom in
the STRONG product.
\end{lemma}
\begin{proof}
Every \(J\)-class is infinite
under SAP: choose a diversity
\(J\)-neighbour of one point
and strongly amalgamate arbitrarily
many copies of this two-point
structure over that point.
The same holds for \(K\).

The STRONG product is contained
in the RAW saturated relation.
Conversely suppose two DISTINCT
real endpoints \(x,y\) have
an old atomic colour in the
RAW saturation. Choose a
witnessing real walk
\(xJu\,t\,vKy\).
If \(x=u\), choose an anchor
\(a\in J[x]\) different from
the finitely many existing
vertices, and strongly amalgamate
two copies of the COMPLETE
old structure on
\(\{a,x,v,y\}\) over the
complete base \(\{a,v,y\}\).
The resulting second copy \(u'\)
of \(x\) is different from
\(x\), belongs to the SAME
\(J\)-class via the common
anchor \(a\), and still has
the old atom \(t\) to \(v\).
Thus \(xJ^\circ u'\,t\,v\).
If \(x\ne u\) simply put
\(u'=u\).

If \(v=y\), choose a new
anchor \(b\in K[y]\) away
from the current finite set,
and strongly duplicate \(y\)
over the complete diagram on
the other existing vertices
and \(b\), fixing \(x,u'\).
The second copy \(v'\ne y\)
still has atom \(t\) to \(u'\)
and belongs to the same
\(K\)-class as \(y\) via \(b\).
Thus \(u'tv'K^\circ y\).
When \(v\ne y\) put \(v'=v\).

We now have a REAL represented
word walk between the original
distinct endpoints with three
DIVERSITY factors
\(J^\circ,\{t\},K^\circ\).
Use SAP injectivisation
\cref{lem:typed-injectivisation}
with \(x,y\) pinned to
separate any remaining
nonconsecutive internal
occurrences. Its endpoint
colour is unchanged and
lies in the STRONG product,
as required.

For the two-factor identity, take
distinct old points \(x,y\) and
an ACTUAL RAW witness \(xJzKy\).
The witness \(z\) may coincide
with one endpoint; it would be
INVALID to try to clone \(z\)
over a base CONTAINING \(z\).
We distinguish all three cases,
using old SAP and the infinite
nonzero equivalence classes
established above.

If \(z\notin\{x,y\}\),
strongly amalgamate two copies
of the COMPLETE old structure
on \(\{x,y,z\}\) over the
COMPLETE pair \(\{x,y\}\).
Keep the first old \(z\),
and call its second copy \(z'\).
Then \(z'\ne x,y\),
\(xJz'\) and \(z'Ky\).

If \(z=x\), so \(xKy\),
choose a fresh point
\(a\in J[x]\setminus\{x,y\}\).
Strongly amalgamate two copies
of the COMPLETE old structure
on \(\{a,x,y\}\) over the
COMPLETE pair \(\{a,y\}\),
retaining \(x\) from the first
copy and writing \(z'\)
for the second copy of \(x\).
The second copy is distinct
from \(x,a,y\); it satisfies
\(z'JaJx\) and \(z'Ky\)
(the original \(xKy\) relation
is preserved to the common \(y\)).
Thus \(xJ^\circ z'K^\circ y\).

If \(z=y\), so \(xJy\),
choose a fresh point
\(b\in K[y]\setminus\{x,y\}\).
Strongly amalgamate two copies
of the COMPLETE old structure
on \(\{x,y,b\}\) over
\(\{x,b\}\), retaining the
first \(y\) and calling its
second copy \(z'\).
Then \(z'\ne x,y,b\),
\(xJz'\) and \(z'KbKy\),
so again \(xJ^\circ z'K^\circ y\).

In the first case the same
STRONG conclusion holds because
\(z'\ne x,y\). Every amalgamation
keeps the ORIGINAL old fine
atomic type of the pinned
endpoints \(x,y\) unchanged.
Hence IN ALL cases there exists
a distinct-middle-point walk
\(xJ^\circ z'K^\circ y\),
with NO equality-labelled
STRONG atom and NO implicit
identification of SAP with AP.
Thus every diversity
output of RAW \(J;K\)
lies in the strong product.
The reverse inclusion
is immediate; commuting
same-sort invariant
equivalences have
\(J;K=J\vee K\).
This proves
\eqref{eq:typed-equivalence-join-strong}.

With an identity-only
endpoint kernel simply omit
the corresponding exterior
step throughout.
\end{proof}

For a fixed ordered pair of sorts \(a,b\), call a non-empty subset of the
atomic \(a\)-\(b\) types \emph{typed exact} if it occurs as a fixed-point
path-consistency domain of a completable finite typed partial structure; write
\(\Dat^{ab}\) for the finite family of these domains.

\begin{lemma}[Typed exact-domain calculus]
\label{lem:typed-exact-domain-calculus}
For every ordered pair of current unary sorts \(a,b\),
\[
 \Dat^{ab}
 =\left\{\bigcap\mathcal F:
   \mathcal F\subseteq\Rch^{ab}\text{ finite and }
   \bigcap\mathcal F\ne\varnothing\right\}.
\tag{TED}\label{eq:typed-exact-domains}
\]
Finite pairwise-intersecting families in one \(\Dat^{ab}\) have non-empty
total intersection.  Moreover, for any sort triple \(a,b,c\), a triple
\[
 D_{ab}\in\Dat^{ab},\qquad
 D_{bc}\in\Dat^{bc},\qquad
 D_{ca}\in\Dat^{ca}
\]
occurs simultaneously as the three fixed-point domains of one completable
typed triangle network if and only if it is support-consistent under the typed
Peircean products.
\end{lemma}

\begin{proof}
This is exactly \cref{lem:typed-PC-fixedpoint-exact}
in the present current named-sorted strong SAP+CCO
source. The canonical typed walk formula
\cref{lem:typed-PC-walk-formula} identifies PC
fixed points with nonempty path-profile meets;
the three-domain assertion is the same typed
Peirce support-consistency condition.
\end{proof}

\begin{lemma}[Original quotient first-hit cover interpolation]
\label{lem:typed-original-full-cover-interpolation}
Let $M$ be a finite-complete homogeneous
binary source with a named transitive
original real sort $a$, finite symmetric
SAME-SORT fine atomic palette, original
visible SAP and semantic HOMOMORPHIC CCO.
Let $\Id_a<J\prec K$ be an ACTUAL FULL
original invariant-equivalence cover,
with $J$ MEET-IRREDUCIBLE (so $K$
is its UNIQUE FULL cover).
Assume that both $J^\circ$ and $K^\circ$
are single positive STRONG OLD typed
word endpoint profiles; their private
word interiors may use arbitrary named
original sorts. Then for EVERY nonempty
positive STRONG typed $a$--$a$ word profile
$P$, we have
\[
 P\cap J^\circ\ne\varnothing,\qquad
 P\setminus J\ne\varnothing
   \quad\Longrightarrow\quad
                   P\cap(K\setminus J)\ne\varnothing.
\tag{OWI}\label{eq:typed-original-full-cover-interpolation}
\]
There is NO exact-domain product calculus,
Set-Pasch, lattice distributivity or hidden
SAP assumption in this original theorem.
\end{lemma}
\begin{proof}
Suppose for contradiction that
$P\cap(K\setminus J)=\varnothing$.
The old SAP full $J$-saturation
\[
           S=J^\circ\starprod P\starprod J^\circ
\]
is positive and is the FULL old diversity
pullback $\Phi(\bar S)$ of one invariant
binary quotient relation $\bar S$ on
$Q=M_a/J$. The quotient relation contains
the TRUE quotient diagonal $\Delta_Q$,
since $P$ meets $J^\circ$ and full
$J$-saturation fills the diagonal
quotient orbital. Saturation cannot
change whether endpoint J-classes
are K-related; since P avoids
$K\setminus J$, we have
\[
                   \bar S\cap(K/J)=\Delta_Q.
\tag{OWI-S}\label{eq:typed-old-interpolation-S}
\]
But $P$ exits $J$ and avoids $K\setminus J$,
so $\bar S$ has an edge OUTSIDE $K/J$.
Its connectedness equivalence H on Q
lifts to an ORIGINAL invariant
equivalence strictly above J.
Since J has UNIQUE FULL cover K,
we have $K\le H$. Thus some quotient
pair in $\bar U=(K/J)\setminus\Delta_Q$
is joined by an $\bar S$-walk.
Because $\Delta_Q\subseteq\bar S$,
the powers of $\bar S$ increase.
Choose the LEAST $n\ge2$ such that
$\bar S^{;n}$ meets $\bar U$.
Put
\[
              A=\bar S^{;(n-1)},\qquad
              T=\bar S;\bar U.
\]
Minimality and \eqref{eq:typed-old-interpolation-S}
give
\[
 A\cap(K/J)=\Delta_Q,\qquad
 T\cap(K/J)=\bar U.
\tag{OWI-ST}\label{eq:typed-old-interpolation-AT}
\]
For the second equality, an $\bar S$-edge
inside a K-component must be quotient
diagonal, whereas an $\bar U$-edge
always stays inside that K-component.
The quotient diagonal of $\bar S$
also supplies every external U edge.

By the choice of n,
$(A;\bar S)\cap\bar U\ne\varnothing$.
Peirce rotation of a witnessed ACTUAL
quotient triangle, using symmetry of
all original same-sort quotient orbitals,
gives
\[
                         A\cap T\ne\varnothing.
\tag{OWI-P}\label{eq:typed-old-interpolation-peirce}
\]

By the original positive cover-layer theorem
\cref{thm:typed-original-full-cover-layer-word},
the quotient layer $\bar U$ has a SINGLE
positive STRONG OLD word pullback U, since
$J^\circ$ is positive. The three OLD
diversity menus
\[
               \Phi(A),\quad \Phi(T),
                           \quad \Phi(K/J)=K^\circ
\]
are therefore each SINGLE positive STRONG
typed OLD word profiles. Indeed by the
exact quotient STRONG composition
\cref{lem:typed-quotient-strong-semigroup},
$\Phi(A)=S^{\starprod(n-1)}$ and
$\Phi(T)=S\starprod U$.
Their pairwise intersections are NONEMPTY:
the first and third meet in $J^\circ$
by \eqref{eq:typed-old-interpolation-AT},
the second and third meet in U,
and the first two meet by
\eqref{eq:typed-old-interpolation-peirce}
and the injectivity of $\Phi$.
Their TOTAL intersection is EMPTY:
inside K/J the quotient supports
of A,T are respectively the
disjoint $\Delta_Q,\bar U$.

This contradicts the original positive
STRONG word 2-Helly consequence of
semantic CCO+visible SAP: private positive
OLD typed paths for pairwise intersecting
word profiles give homomorphically
realisable specified cycles, a global
CCO map, and then fully DISTINCT roots
by SAP injectivisation
\cref{lem:typed-injectivisation}.
Hence P meets $K\setminus J$.
\end{proof}

\begin{lemma}[Typed cover interpolation]
\label{lem:typed-cover-interpolation}
Let \(J\) be a represented proper meet-irreducible on a real sort \(a\), let
\(K=J^+\) be its unique cover, and let \(P\in\Rch^{aa}\), where the path
realising \(P\) may pass through arbitrary intermediate sorts.  If
\[
 P\cap J\ne\varnothing,
 \qquad
 P\setminus J\ne\varnothing,
\]
then \(P\cap(K\setminus J)\ne\varnothing\).
\end{lemma}

\begin{proof}
A SHORT independent proof is now available
directly from original quotient geometry.
The $J=\Id_a$ case is VACUOUS because a
STRONG diversity profile contains no
logical real identity. Otherwise the
TERMINAL full NR2 hypothesis and
\cref{lem:typed-all-invariant-equivalences-raw-word}
make the represented exact MI J and
its exact cover K an actual FULL
invariant MI and its UNIQUE FULL cover.
Both nonzero full RAW equivalences J,K
have positive STRONG word profiles by
original visible SAP injectivisation.
The quotient first-hit theorem
\cref{lem:typed-original-full-cover-interpolation}
therefore applies DIRECTLY to the
given reachable profile P, including
when P has intermediate vertices of
other original named sorts. This
proves the asserted cover interpolation
without the older Set-Pasch,
distributivity or minimal-power descent.

\end{proof}

\begin{lemma}[Original STRONG exact-domain MI-primality from one cover]
\label{lem:typed-original-strong-exact-prime}
In the setting of
\cref{lem:typed-original-full-cover-interpolation},
let $J>\Id_a$ be the FULL original
meet-irreducible with unique FULL
cover $K=J^+$, and let $D_1,D_2$
be NONEMPTY finite intersections
of genuinely positive OLD STRONG
typed $a$--$a$ word profiles.
(The empty intersection means
the full diversity palette and
is handled separately.) If
\[
     D_1\cap D_2\ne\varnothing,\qquad
     D_1\cap D_2\subseteq J^\circ,
\]
then
\[
                 D_1\subseteq J^\circ
                   \quad\text{or}\quad
                 D_2\subseteq J^\circ.
\tag{OWP}\label{eq:typed-original-strong-exact-prime}
\]
The proof uses only ORIGINAL
SAP+homomorphic CCO positive-word
2-Helly, the direct quotient
cover-layer positivity, and
the first-hit quotient
interpolation. It does NOT
assume a typed exact-domain
product formula, distributivity,
original Set-Pasch, or hidden SAP.
\end{lemma}
\begin{proof}
If one $D_i$ is the full
diversity palette (empty family
of defining words), the premise
already makes the OTHER
$D_j\subseteq J^\circ$;
so assume nonempty defining
families $\mathcal P_1,\mathcal P_2$.
Suppose both $D_i$ exit $J$.
Fix one old fine diversity atom
$c\in D_1\cap D_2\subseteq J^\circ$.
For EVERY defining positive
profile $P\in\mathcal P_i$,
we have $c\in P$ and
$P\setminus J\ne\varnothing$,
because $D_i\subseteq P$
and $D_i$ exits J.
The direct first-hit interpolation
\cref{lem:typed-original-full-cover-interpolation}
therefore shows that EVERY such
P meets the POSITIVE external
layer $U=K\setminus J$.
This layer is one old STRONG
positive word profile by
\cref{thm:typed-original-full-cover-layer-word}.

For each fixed i, the finite
family $\{U\}\cup\mathcal P_i$
is PAIRWISE-intersecting: U
meets every P by interpolation,
and all P meet each other at
the SAME original fine atom c.
The original positive-word
2-Helly theorem (private paths,
homomorphic CCO and original SAP
injectivisation) thus gives
$D_i\cap U\ne\varnothing$ for
both i.

Finally $D_1,D_2,U$ are
pairwise-intersecting finite
intersections of OLD positive
STRONG word profiles. Flatten
their defining families.
Every defining word from one
$D_i$ meets every defining
word from the other through
the assumed nonempty
$D_1\cap D_2$; similarly the
$D_i$--U pairs meet by the
previous paragraph. Positive
word 2-Helly yields
$D_1\cap D_2\cap U
\ne\varnothing$, contradicting
$D_1\cap D_2\subseteq J^\circ$
and $J^\circ\cap U=\varnothing$.
This proves (OWP).
\end{proof}

\begin{lemma}[Typed prime intersection]
\label{lem:typed-prime-intersection}
Let \(J\) be a represented proper meet-irreducible on sort \(a\).  If
\(D_1,D_2\in\Dat^{aa}\) have non-empty intersection and
\[
 D_1\cap D_2\subseteq J,
\]
then \(D_1\subseteq J\) or \(D_2\subseteq J\).
\end{lemma}

\begin{proof}
The $\Id_a$ case is vacuous because
$D_1\cap D_2$ is a NONEMPTY
STRONG diversity domain, never
the logical equality atom.
Otherwise the terminal ORIGINAL
full NR2 and the CCO-free
\cref{lem:typed-all-invariant-equivalences-raw-word}
identify every original invariant
equivalence with an ACTUAL FULL
represented exact equivalence.
Thus J is an ORIGINAL FULL MI
with its unique FULL invariant
cover K. Original visible SAP
injectivises positive RAW words
for the nonzero J,K diversity
palettes. By
\cref{lem:typed-exact-domain-calculus},
each $D_i$ is a finite nonempty
intersection of positive STRONG
typed word profiles, with an
empty defining family handled
trivially as FULL diversity top.
The ORIGINAL strong exact-prime
\cref{lem:typed-original-strong-exact-prime}
therefore applies DIRECTLY and
proves the conclusion.
This uses no extra hidden SAP
and avoids importing the
historical Set-Pasch/descent
proof of cover interpolation.

\end{proof}

\noindent\emph{Foundational transport imported.}
The general invariant typed quotient-transport statement
\cref{lem:typed-quotient-transport} and its proof now occur in
Part~I, \cref{sec:typed-simultaneous-exact}.
Here it is applied to the already named original real sorts;
no quotient-section, exactness of a transported nonexact relation,
or SAP is inferred merely from that lemma.

\begin{lemma}[No algebraic real quotient under SAP]
\label{lem:no-real-algebraic-transition}
Assume visible SAP.  Let \(R\in\Rel^{ab}\) and let \(E\) be a proper
NONIDENTITY definable equivalence on sort \(a\), so \(E>\Id_a\), with
\(RR^{\smile}\subseteq E\).  Then the transported equivalence \(E^R\) on
sort \(b\) is proper; in particular \(E^R\) is not equality.
Moreover, if \(R\in\Rel^{ab}\) satisfies
\(RR^{\smile}=1'_a\), then \(a=b\) and \(R=1'_a\).
\end{lemma}

\begin{proof}
If \(E^R\) were equality, then for every \(x\in a\) there would be exactly
one \(y\in b\) with \(xRy\).  Indeed existence follows from totality of
\(R\), while two such points \(y,y'\) satisfy
\(yE^Ry'\) because \(xEx\), and hence \(y=y'\).  Thus \(R\) is the graph of
a total \(0\)-definable function \(f:a\to b\), constant on \(E\)-classes.
Choose \(x\) and put \(y=f(x)\).  Unless \(a=b\) and \(y=x\), strongly
amalgamating two copies of the finite two-point structure generated by
\(x,y\) over the one-point substructure \(\{x\}\) would have to keep the two
copies of \(y\) distinct, while functionality of \(R\) forces them to be the
same point.  This contradicts SAP.  Hence every value of a \(0\)-definable
total function between real sorts is the source point; by transitivity of the
current unary sorts this is possible only when source and target sorts agree
and the function is the identity.  But for
\(R=1'_a\) we have \(E^R=E\), contradicting properness of \(E\) if
\(E^R\) were equality.  This proves the first assertion.

For the FINAL assertion the direction of functionality is
the CONVERSE one. If \(RR^\smile=1'_a\),
then every \(y\) in sort \(b\) has exactly ONE
\(R\)-predecessor \(x\) in sort \(a\):
existence follows from surjectivity of
the nonempty invariant \(R\), while
two predecessors \(x,x'\) satisfy
\(xRR^\smile x'\) and hence \(x=x'\).
Thus \(R^\smile\) is the graph of
a total \(0\)-definable function \(b\to a\).
Visible SAP, by the same two-copy
strong amalgam over the INPUT point
\(y\), forces this function to be
literal old identity. Consequently
\(a=b\) and \(R=1'_a\).
No equality of \(E^R\) was asserted
or needed for this converse-direction argument.
\end{proof}

\begin{lemma}[Identity-containing quotient transport is redundant]
\label{lem:typed-identity-transport-redundant}
Let \(R\in\Rel^{aa}\) contain the ACTUAL
logical diagonal \(\Id_a\), and let
\(J\) be a represented proper same-sort
equivalence with \(R;R^\smile\subseteq J\).
Then
\[
 R\subseteq J,\qquad
 J^R=R^\smile;J;R=J,\qquad
 q_{R,J}=\operatorname{id}_{a/J}.
\tag{QT-Id}\label{eq:typed-identity-transport}
\]
Such an \(R\) gives no new unary
quotient-transition function. Its map
is the existing identity on
\(B_{a,J}\), without any artificial
STRONG identity-labelled edge.
This includes \(R=\Id_a\)
and any mixture of identity and
genuine diversity atoms.
\end{lemma}
\begin{proof}
Same-sort old orbitals are symmetric,
so \(R^\smile=R\). Since
\(\Id_a\subseteq R\),
\[
 R=\Id_a;R\subseteq R;R^\smile
                      \subseteq J.
\]
Thus \(R^\smile;J;R\subseteq
J;J;J=J\); the reverse inclusion
follows from \(\Id_a\subseteq R\),
so \(J^R=J\). For any \(xRy\)
we have \(xJy\), hence the
induced quotient map sends
each class to itself. Its
graph is the existing unary
identity term, not an additional
positive-path constraint.
\end{proof}

\begin{lemma}[Finite typed quotient-transition composition normal form]
\label{lem:typed-transition-composition-normal-form}
Let \(R\in\Rel^{ab}\),
\(S\in\Rel^{bc}\) be
nonempty ACTUAL typed
invariant relations.
Let \(E\) be a represented
proper equivalence on
sort \(a\), and assume
\[
 RR^\smile\subseteq E,\qquad
 F=E^R,\qquad
 SS^\smile\subseteq F .
\]
Put \(T=R;S\) as an
equality-inclusive RAW
represented relation.
Then
\[
 TT^\smile\subseteq E,\qquad
 E^T=F^S,\qquad
 q_{T,E}=q_{S,F}\circ q_{R,E}.
\tag{TQC}\label{eq:typed-transition-composition}
\]
For any represented
equivalence \(H\ge F\)
on \(b\), let
\(H_0=(q_{R,E})^{-1}(H)\)
be its unique pullback
in the interval
\([E,\Top_a]\). Then
\(H_0^R=H\) and
\[
 u_{F,H}\circ q_{R,E}
 =q_{R,H_0}\circ u_{E,H_0}.
\tag{TQU}\label{eq:typed-upward-interchange}
\]
If \(H\) is a represented
proper meet-irreducible,
so is \(H_0\). If
\(a=c\) and the ACTUAL
diagonal \(\Id_a\)
belongs to \(T\), then
the composite quotient
map in (TQC) is exactly
the IDENTITY on \(a/E\),
by \cref{lem:typed-identity-transport-redundant};
it introduces no synthetic
STRONG diversity identity
or additional named unary
transition.

Consequently every
transition-only unary
composite between
represented proper
quotient sorts equals
either one quotient
transition \(q_{T,E}\)
of an ACTUAL member of
one of the finitely many
families \(\Rel^{ac}\),
or the old unary identity.
After moving upward maps
leftwards by (TQU),
every unary coordinate
term has one initial
projection/upward map
and one such transition
(or identity). Hence
there are finitely many
unary term FUNCTIONS
from each hidden sort.
\end{lemma}
\begin{proof}
The relation \(T=R;S\)
is nonempty and total and
surjective between its
endpoint sorts because
\(R,S\) have these
properties on transitive
original named sorts.
If \(xTz\) and \(x'Tz\),
choose intermediate
\(y,y'\) with
\(xRySz\) and
\(x'Ry'Sz\).
Since \(ySS^\smile y'\),
we have \(yFy'\).
By the bijection
\(q_{R,E}:a/E\to b/F\)
of
\cref{lem:typed-quotient-transport},
the corresponding
\(E\)-classes of \(x,x'\)
are equal. Thus
\(TT^\smile\subseteq E\).
For a target point \(z\),
the unique source
\(E\)-class associated
through \(T\) is the
one obtained by applying
the inverse quotient
bijections for \(S\)
and then \(R\).
Hence the target kernel
is exactly \(F^S\), and
the forward quotient
map is their composition,
proving (TQC).

The interval-lattice
isomorphism in
\cref{lem:typed-quotient-transport}
sends \(H_0\) to \(H\).
By the definition of
an upward quotient map,
the two compositions in
(TQU) send any
\(E\)-class \(q\)
to the same \(H\)-class
of its \(R\)-neighbours.
They therefore agree
pointwise. A proper
meet-irreducible in
\([F,\Top_b]\) has a
unique cover there and
hence the pullback has
a unique cover in
\([E,\Top_a]\), which is
also its unique cover
in the full actual
equivalence lattice.
This proves the
meet-irreducibility
assertion.

When \(T\) contains
logical identity the
identity-redundancy lemma
directly yields
\(T\subseteq E\) and
\(q_{T,E}=\mathrm{id}\).
Do NOT infer the
existence of a
diversity identity atom
from this RAW conclusion.

Repeated use of (TQC)
replaces any finite
transition-only composite
by the one relation
product of its original
typed \(R\)-relations.
Such a product is
again a member of
\(\Rel^{ac}\), since it
is actual invariant
and is a union of
finitely many original
typed atomic orbitals.
If it contains identity,
the preceding paragraph
replaces its quotient map
by an existing unary
identity. Repeated use
of (TQU) moves every
upward quotient map
toward the source.
There are finitely many
named real sorts, proper
represented quotient sorts
and invariant binary
relations \(\Rel^{ac}\),
so only finitely many
distinct normal-form
unary functions can occur.
\end{proof}

Let \(\mathcal I_a\) be the proper represented meet-irreducibles on sort
\(a\).  Equality is represented by the real sort itself; by
\cref{lem:no-real-algebraic-transition} no non-trivial quotient transition
can enter a real sort.  Thus we introduce hidden coordinate sorts
\(B_{a,J}\) only for \(J\in\mathcal I_a\), together with the old projections
from the real sort and upward quotient maps between represented coordinates.
Whenever \(J\in\mathcal I_a\), \(R\in\Rel^{ab}\),
\(RR^{\smile}\subseteq J\), and the ACTUAL logical
identity is NOT contained in \(R\), introduce
the unary quotient-transition map
\[
 \pi_{R,J}:B_{a,J}\longrightarrow B_{b,J^R},
\tag{QT3}\label{eq:typed-transition-function}
\]
where \(J^R\in\mathcal I_b\) follows from the FULL invariant
interval transport, the SEPARATE terminal RAW exactness theorem,
and the SAP noncollapse statement, respectively:
\cref{lem:typed-quotient-transport,lem:typed-terminal-all-equivalences-exact,lem:no-real-algebraic-transition}.
The generic transport lemma alone does NOT make an arbitrary
transported invariant equivalence FULL RAW exact.
Quotient transitions are bijections, so include their inverses and all
composition identities supplied by \cref{lem:typed-quotient-transport}.
If a relational composite contains genuine identity,
\cref{lem:typed-identity-transport-redundant}
reduces it to an existing quotient identity term;
no synthetic strong identity edge is introduced.
For
the cylinder notation it is convenient to regard the real sort \(a\) as the
equality coordinate \(a/1'_a\).  Put
\[
 \overline{\mathcal I}_a=\mathcal I_a\cup\{1'_a\},
 \qquad B_{a,1'_a}:=a,
\]
where a \(1'_a\)-class is a singleton.
For every nonzero represented proper
same-sort equivalence \(J\)
used as a STRONG factor below,
its diversity palette \(J^\circ\)
is nonempty and itself a
POSITIVE typed reachable profile:
\cref{lem:typed-terminal-all-equivalences-exact}
represents the FULL \(J\)
by one equality-inclusive
old atomic word in this
TERMINAL typed source, and
\cref{lem:typed-injectivisation}
identifies its distinct-endpoint
profile with \(J^\circ\).
That terminal reachability
lemma uses
\cref{lem:typed-terminal-pure-power}
and the original RAW exact
calculus, not the cylinder
construction here; the forward
reference does NOT assume
quotient-transition closure.
For
\(J\in\overline{\mathcal I}_a\),
\(K\in\overline{\mathcal I}_b\), and an atomic \(a\)-\(b\) pair type \(t\),
define the representative cylinder by omitting identity
saturation and using the explicit diversity
parts of nonzero equivalences from
\cref{lem:typed-raw-strong-saturation}:
\[
 \operatorname{Cyl}_{J,K}(t)=
 \begin{cases}
 \{t\},&J=1'_a,\ K=1'_b,\\
 \{t\}\starprod K^\circ,&J=1'_a,\ K\neq1'_b,\\
 J^\circ\starprod\{t\},&J\neq1'_a,\ K=1'_b,\\
 J^\circ\starprod\{t\}\starprod K^\circ,&J\neq1'_a,\ K\neq1'_b.
 \end{cases}
\tag{QT-Cyl}\label{eq:typed-cylinder}
\]

\begin{lemma}[Typed cylinder realisation and partition]
\label{lem:typed-cylinder-realisation}
Let \(J\in\overline{\mathcal I}_a\) and
\(K\in\overline{\mathcal I}_b\).  Let \(q\in a/J\) and \(r\in b/K\),
where an equality-coordinate class is identified with its unique real point.
Suppose distinct representatives \(x\in q\), \(y\in r\) have atomic pair type \(t\).
Then the set of all atomic pair types realised by choices of representatives
\(x'\in q\), \(y'\in r\) with \(x'\neq y'\) is exactly
\[
 \operatorname{Cyl}_{J,K}(t).
\]
Consequently this cylinder depends only on the ordered pair \((q,r)\), not
on the chosen representatives.  For every compatible ordered pair of
represented coordinate positions, including a real-sort endpoint interpreted
as equality, the distinct cylinders therefore partition the ordered pairs of
\emph{distinct} corresponding vertices (with the diagonal omitted when the
two endpoint positions are the same hidden sort); when both endpoints are real
this is simply the old atomic pair-type partition off the diagonal.
\end{lemma}

\begin{proof}
If \(x'\in q\) and \(y'\in r\), then the left exterior step has type in
\(J\) when \(J\neq1'_a\) and is absent when \(J=1'_a\); similarly the
right exterior step has type in \(K\) unless \(K=1'_b\).  Together with the
middle type \(t\), this gives exactly the appropriate case of
\eqref{eq:typed-cylinder} for the type of \(x'y'\).

Conversely, let \(s\in\operatorname{Cyl}_{J,K}(t)\).  By the semantic
definition of the typed Peircean product, realise the corresponding path in
the homogeneous limit: its middle edge has type \(t\), with a left exterior
edge in \(J\) precisely when \(J\neq1'_a\) and a right exterior edge in
\(K\) precisely when \(K\neq1'_b\).  Thus the path has between two and four
vertices.  An automorphism sends its middle edge to \((x,y)\).  Every exterior
vertex is then sent into the required quotient class, while an
equality-coordinate endpoint is already the corresponding middle vertex.
The endpoint type remains \(s\).  Thus every member of the displayed cylinder
is realised by distinct representatives of \((q,r)\).  The final assertion follows by
choosing representatives of each coordinate pair and applying the first
assertion twice.
\end{proof}

Accordingly, between every compatible ordered pair of represented coordinate
positions with at least one proper quotient coordinate we name one binary
semantic-cylinder relation for each distinct cylinder in
\eqref{eq:typed-cylinder}, interpreted by the corresponding part of the
partition in \cref{lem:typed-cylinder-realisation}.  Real--real cylinders are
already named by the old atomic pair-type predicates.

\begin{lemma}[Quotient-transition cylinders]
\label{lem:typed-transition-cylinder}
Let \(J\in\mathcal I_a\), let \(R\in\Rel^{ab}\)
be nonempty and contain NO ACTUAL logical identity,
assume \(RR^{\smile}\subseteq J\), and put \(K=J^R\).
For every DIVERSITY atomic pair type
\(t\in R\),
\[
 \operatorname{Cyl}_{J,K}(t)
   =J^\circ\starprod R\starprod K^\circ.
\tag{QT-Graph}\label{eq:typed-transition-cylinder}
\]
In particular the left-hand side is independent of \(t\).  If
\(q\in a/J\) and \(r\in b/K\), then
\[
 \pi_{R,J}(q)=r
 \quad\Longleftrightarrow\quad
 \text{the type of representatives of \((q,r)\) lies in
 \(J^\circ\starprod R\starprod K^\circ\)}.
\]
Thus the graph of every quotient-transition function is represented by one
reachable visible pair profile.
\end{lemma}

\begin{proof}
For every old diversity atomic $t\in R$,
we have $tt^\smile\subseteq RR^\smile
\subseteq J$. The atom $t$ is total and
surjective between the original transitive
sorts. For every $y\in b$, choose
$x\,t\,y$; as $t\subseteq R$, the unique
source $J$-class is the same for $t$ and
for $R$. Consequently
\[
 \lambda_{t,J}=\lambda_{R,J},
 \qquad J^t=J^R=K,\qquad
 q_{t,J}=q_{R,J}.
\]
For ANY matched quotient pair
$r=q_{R,J}(q)$, each $x\in q$ has a
$t$-neighbour $y\in r$. Thus every
graph pair admits one actual old
$t$-representative pair. By original
two-point homogeneity, ALL these pairs
belong to ONE complete $\Aut(M)$-orbital
of quotient pairs, including when $a\ne b$.
This orbit is the graph of $q_{R,J}$.

The complete OLD diversity spectrum of
one matched pair is
$(J;\{t\};K)^\circ$ by
\cref{lem:typed-cylinder-realisation}.
It is independent of the chosen atomic
$t\in R$, since each such atom induces
the SAME quotient graph. Taking their
union therefore gives
\[
 (J;R;K)^\circ
   =(J;\{t\};K)^\circ
   =\operatorname{Cyl}_{J,K}(t).
\]
Here all semicolons are actual RAW
equality-inclusive relational products.
Both kernels are NONZERO under visible
SAP, so \cref{lem:typed-raw-strong-saturation}
converts this diversity slice into
$J^\circ\starprod R\starprod K^\circ$.
It equals the corresponding product
with the SINGLE atom $t$ and is
STRONG reachable by the terminal
positive word representation of
$J^\circ,K^\circ$. This proves
\eqref{eq:typed-transition-cylinder}.

Conversely, if a distinct OLD
representative pair has a fine atom
in that graph spectrum, original
two-point homogeneity carries it
to a representative pair of a matched
quotient class pair. Thus its classes
satisfy $r=q_{R,J}(q)$. In a matched
pair EVERY distinct representative
pair has fine atom in the same full
spectrum. Distinct representatives
always exist because proper nonzero
J,K-classes are infinite by SAP.
The possible coincidence of the
HIDDEN quotient vertices is checked
by unary equality and is not an
artificial off-diagonal binary atom.
\end{proof}

\begin{lemma}[Atomic naming of a quotient-functional strong branch]
\label{lem:typed-branch-atomic-transport}
Let \(P\in\Rch^{ab}\) be a nonempty
IDENTITY-FREE positive typed strong
word profile, and let \(J\) be a
represented proper equivalence on
the source sort \(a\) with
\(P;P^\smile\subseteq J\).
Choose ANY old diversity
atomic type \(t\subseteq P\).
Then
\[
      J^t=J^P,\qquad
      \pi_{t,J}=\pi_{P,J}.
\tag{QT-Atom}\label{eq:typed-atomic-transport}
\]
In particular, when \(J\) is a
represented proper meet-irreducible,
the transported target kernel
\(J^P\) is also a represented proper
meet-irreducible under visible SAP,
and the quotient map induced by
the WHOLE strong branch \(P\) is
ALREADY a named unary quotient
transition of the finite hidden
language. The assertion does NOT
identify old fine atomic relations
\(P\) and \(t\); only their
induced quotient maps agree.
\end{lemma}
\begin{proof}
The actual relation \(P\) is
nonempty and invariant between
the transitive named sorts
\(a,b\), so is total and
surjective. The same holds for
the atomic orbital \(t\).
Since \(t\subseteq P\),
\(t;t^\smile\subseteq
P;P^\smile\subseteq J\),
and both quotient transports
are defined by
\cref{lem:typed-quotient-transport}.

For any target point \(y\)
choose an OLD real point
\(x\) with \(x\,t\,y\),
possible by surjectivity.
Then \(x\,P\,y\) too,
so the unique source
\(J\)-class associated to
\(y\) through the
relation \(t\) is the SAME
one associated through \(P\).
Thus their well-defined
maps \(\lambda_{t,J}\)
and \(\lambda_{P,J}\)
coincide POINTWISE on
the actual target sort \(b\).
Their kernel equivalences
\(t^\smile Jt\) and
\(P^\smile JP\) are equal;
the quotient bijections
defined as inverses of the
same source-class labelling
are identical. This gives
\eqref{eq:typed-atomic-transport}.
If \(J\) is proper MI,
transport preserves
meet-irreducibility, and
\cref{lem:no-real-algebraic-transition}
excludes an equality target
under visible SAP. Since
\(t\) is a genuine old
diversity atom, the
symbol \(\pi_{t,J}\)
is one of the explicitly
named IDENTITY-FREE
transition symbols.
\end{proof}

There are only finitely many new symbols.  In particular, functional
transition closure never creates a new visible real point.

